\section{Introduction}
Dark matter halo masses connect galaxy populations to their gravitational environments. A galaxy system provides a variable number of tracers with correlated positions, velocities, and intrinsic properties. Graph neural networks (GNNs) offer a natural representation of these data: galaxies become nodes, relationships become edges, and permutation-invariant pooling produces a halo-level estimate. Existing work already demonstrates GNN halo-mass inference across simulation suites \citep{villanueva2022,garuda2024}. The question here is narrower: how much predictive information survives when interactions in an existing predictor are removed under explicit structural constraints?

Sparsification has several possible purposes, including reducing computation, identifying predictive relationships, or limiting communication. These purposes require different evidence. Preserving predictions does not establish a physical explanation, and removing edges does not establish a speedup when a selector itself requires a full model pass. A scientific evaluation must therefore measure accuracy, selection feasibility, and total inference cost together, against competitive selectors with identical final budgets.

Cosmic-Net combines an edge-conditioned GNN with a learned pair selector and a relative virial-inspired reward. We report its completed A--D campaign as a controlled empirical study of that implemented system. The campaign provides a useful test bed for exact budget accounting and constrained evaluation, despite unresolved provenance in its historical catalog and frozen checkpoint. We retain the explicit research label \code{legacy\_total\_radius}; the measurements are conditional results on that artifact, rather than a certified measurement of a particular spherical-overdensity mass definition.

The study has three aims. First, we describe the actual selection and reward implementation, including the difference between repaired training actions and exact-budget inference. Second, we compare seven selectors using the same frozen predictor, halo population, and physical-pair decoder. Third, we quantify uncertainty in small observed differences and account for selection overhead. The main result is a qualified negative one: at 40\% nominal retention the learned policy is less accurate than the full graph, its point advantage over gradient~$\times$~input is unresolved, and the implemented inference path is substantially slower. These observations motivate concrete changes in evaluation and deployment design without asserting that learned sparsification cannot be useful elsewhere.
